\clearpage
\SourceLeaf{172}\centerline{[ON THE THEORY OF PARALLEL LINES.]}

\bigskip
\hrule
\bigskip

\centerline{[1.]}

\smallskip

\centerline{\textsc{Gauss} to \textsc{Olbers}. \ Braunschweig, 30 July 1806.}

\smallskip

.\ .\ .\ .\ .\ Mr. \textsc{von Zach} also writes to me that you offered to review \textsc{Legendre}'s work on cometary orbits. I shall gladly send on to you in Bremen the copy Mr. \textsc{von Zach} sent me, though I trust you will let me keep it a few weeks longer. On a first look through it, it seems to me to contain much that is beautiful. I find in this book much of what was distinctive in my own method, especially in its earliest form. It seems to be my fate to compete with \textsc{Legendre} in nearly all my theoretical work. So it has been in higher arithmetic, in investigations of transcendental functions connected with the rectification of the ellipse, in the first principles of geometry, and now here again. Thus, e.\,g., the principle I have used since 1794, that to represent as well as possible several quantities that cannot all be represented exactly one must make the sum of the squares a minimum, is also used in \textsc{Legendre}'s work and carried out quite ably.\ .\ .\ .\ .\ .\ .

\bigskip
\hrule
\bigskip

\centerline{[2.]}

\smallskip

\centerline{[\emph{From} \textsc{Schumacher}'s \emph{diary: »Gaussiana«. Göttingen, November 1808.}]}

\smallskip

\{\textsc{Gauss} has reduced the theory of parallel lines to the conclusion that, if the assumed theory were false, there would have to be a constant line whose length was given beforehand, which is absurd. He himself still regards this work as insufficient.\}

\bigskip
\hrule
